\section{Conclusion}
\label{sec:conclusion}

We derive a two-boundary renewal estimator for deciding when a body-fixed lookahead maneuver is worth its schedule cost, then reduce it to a monotone break-even heuristic for lightweight onboard gating.
The estimator separates the physical slew prior from a schedule-dependent break-even spacing, and the optimal-timing result (Proposition~\ref{prop:timing}) collapses the one-dimensional timing search to a closed-form initiation time that incurs zero opportunity cost up to a bounded perturbation.
Together, the break-even margin, timing rule, and pitch search produce an onboard policy that asks a direct engineering question: whether any feasible pointing angle can generate enough expected clear-image equivalents to beat the local replacement burden.

\paragraph{Limitations.}
The policy's value estimator assumes equal task utility (A4), i.i.d.\ cloud states (A3), and constant gap spacing (A1).
Relaxing these---particularly for heterogeneous Pareto-distributed task values---is necessary for realistic commercial operations where a small fraction of high-value tasks dominates schedule utility.
The analysis also assumes perfect lookahead classification once an access is visible.
False positives would insert cloudy tasks into the repaired schedule, while false negatives would suppress useful alternate accesses; quantifying sensitivity to classifier precision and recall is therefore necessary before mapping the method to a specific onboard vision pipeline.

\paragraph{Future work.}
Extending the chain-length analysis to utility-weighted chains would address the Pareto case.
Multi-step lookaheads within a single gap could improve coverage for narrow-FoV instruments.
Constellation extensions with shared belief states across spacecraft could reduce redundant observations, and pre-planning lookahead opportunities in the ground schedule could yield global improvements at the cost of additional computational complexity.
